\section{The corrected price game}\label{sec:price}

This section solves the affected \(UU\) component globally.  The central distinction is between a stationary point on one active-set branch and a best response over the full clipped demand correspondence.

\subsection{Piecewise demand}

For \(0<z<1/3\), aggregate switching demand is
\[
F_z(x)=
\begin{cases}
0,&x\le -z,\\
x+z,&-z\le x\le0,\\
3x+z,&0\le x\le z,\\
4x,&z\le x\le1-z,\\
3x+1-z,&1-z\le x\le1,\\
x+3-z,&1\le x\le1+z,\\
4,&x\ge1+z.
\end{cases}
\]
On a smooth branch \(F_z(x)=ax+b\), simultaneous interior first-order conditions from \eqref{eq:segmentpayoffs} imply
\begin{equation}
p=\frac{8-b}{3a},
\qquad
q=\frac{4+b}{3a},
\qquad
x=\frac{2(2-b)}{3a}.
\label{eq:smoothcandidate}
\end{equation}
Imposing the branch restrictions eliminates every smooth mutual candidate except
\[
p_0=\frac23,\qquad q_0=\frac13,\qquad x_0=\frac13.
\]
For example, the candidate generated by the branch \(0\le x\le z\) satisfies
\[
x-z=\frac{4-11z}{9}>0
\]
throughout \(z<1/3\), so it lies outside that branch.  The symmetric upper branches fail analogously.

Kinks do not supply missing mutual optima.  At the lower slope-increase kinks the poacher's one-sided derivatives have the wrong ordering for a local maximum, while at the upper slope-decrease kinks the incumbent's one-sided derivatives have the wrong ordering.  Flat zero-demand/full-capture regions and zero prices also fail the mutual-best-response conditions.  These observations exhaust the pure candidates.

\subsection{The exact validity threshold}

At \(p_0=2/3\), the incumbent remains globally optimal for all \(z<1/3\).  The poacher, however, has a competing three-active-group vertex
\[
q_H=\frac{2+z}{6}.
\]
Its normalized profit is
\[
\Pi_H=\frac{(2+z)^2}{12},
\]
whereas the all-four-active source price yields
\[
\Pi_0=\frac49.
\]
The gain is therefore
\begin{equation}
G(z)=\Pi_H-\Pi_0
=\frac{3z^2+12z-4}{36}.
\label{eq:gain}
\end{equation}
The positive root inside the economically relevant range is
\[
z_c=\frac{-6+4\sqrt3}{3}.
\]

\begin{proposition}[Pure-price correspondence]\label{prop:pure}
In the baseline \(UU\) segment game with nonnegative prices and \(0<z<1/3\):
\begin{enumerate}
\item if \(0<z\le z_c\), \((p,q)=(2/3,1/3)\) is the unique pure-strategy Nash equilibrium;
\item if \(z_c<z<1/3\), no pure-strategy Nash equilibrium exists.
\end{enumerate}
At \(z=z_c\), the poacher has a second payoff-equal best reply to \(p=2/3\), but that action does not form a mutual best response with the incumbent.
\end{proposition}

\begin{proof}
The branch calculation \eqref{eq:smoothcandidate}, the kink one-sided derivative tests, and the boundary tests leave only \((2/3,1/3)\) as a possible pure equilibrium.  At that point the incumbent is globally best responding.  Equation~\eqref{eq:gain} shows that the poacher is globally best responding at \(q=1/3\) if and only if \(z\le z_c\).  For \(z>z_c\) the remaining three-active vertex is strictly better, and because all other mutual pure candidates have already been excluded, no pure equilibrium remains.  At equality the second poaching reply changes the incumbent's best response, so it does not create another equilibrium.
\end{proof}

The exact regression case used throughout the verification suite takes \((\sigma,\Delta)=(25,8)\).  Then \(z=8/25>z_c\), \(q_0=25/3\), \(q_H=29/3\), and the dimensioned gain is
\[
25G(8/25)=\frac{23}{225}>0.
\]
This is a finite feasible deviation inside the source restriction, not a limiting or off-domain example.

\subsection{A certified mixed continuation}

The absence of a pure equilibrium raises a continuation problem for the earlier stages.  We construct one explicit equilibrium rather than impose an ad hoc selection rule.

For \(z_c<z<1/3\), define
\begin{equation}
p^*=z\frac{3+2\sqrt3}{3},
\quad
q_L=z\frac{3+2\sqrt3}{6},
\quad
q_H=z\frac{2+\sqrt3}{3},
\label{eq:mixedprices}
\end{equation}
and
\begin{equation}
\lambda(z)
=\frac{4-(5+3\sqrt3)z}{z(1+\sqrt3)}.
\label{eq:lambda}
\end{equation}

\begin{proposition}[Selected two-point continuation]\label{prop:mixed}
For every \(z_c<z<1/3\), the following is a Nash equilibrium of a baseline \(UU\) segment: the incumbent sets \(p^*\) with probability one; the poacher sets \(q_L\) with probability \(\lambda(z)\) and \(q_H\) with probability \(1-\lambda(z)\).
\end{proposition}

The two poaching atoms are globally optimal against \(p^*\) and deliver identical profit.  The mixing probability in \eqref{eq:lambda} makes \(p^*\) a global incumbent best response.  Direct comparison of every smooth branch and kink verifies globality, and a separately written primitive-demand optimizer reproduces the best responses without using the symbolic branch labels.

The resulting expected normalized component profits are
\begin{equation}
E\pi_I=
\frac{z(5+3\sqrt3)(2-z)}{3},
\qquad
E\pi_P=
\frac{z^2(7+4\sqrt3)}{3}.
\label{eq:mixedcomponentprofit}
\end{equation}
A firm's no-information profit, summing its incumbent and poaching roles, is
\begin{equation}
M(z)=
\frac{z[(2+\sqrt3)z+10+6\sqrt3]}{3}.
\label{eq:M}
\end{equation}

Proposition~\ref{prop:mixed} is an existence result.  We have searched for nearby alternative two-atom equilibria and found no numerically distinct low-regret candidate, but that falsification search is not a proof of unrestricted mixed-equilibrium uniqueness.  Hence the remainder of the high-\(z\) backward induction is conditional on the equilibrium in Proposition~\ref{prop:mixed}.

\subsection{Other terminal subgames}

The failure is localized to the common-control \(UU\) component.  Reconstructing prices from clipped type-level demand shows:
\begin{itemize}
\item \(TT+TT\), used under full information, remains globally valid;
\item \(TT+UU\), used when one firm invests and information is shared, inherits the corrected \(UU\) continuation;
\item \(TU+TU\), used when both firms invest without sharing, remains globally valid;
\item \(TU+UU\), used when one firm invests without sharing, inherits the corrected \(UU\) continuation;
\item \(TT+TU\), used in the asymmetric sharing/both-invest history, remains globally valid.
\end{itemize}
The correction therefore enters backward induction only through histories containing a \(UU\) component.
